\documentclass[10pt,aspectratio=169]{beamer}
\usepackage[utf8]{inputenc}
\usepackage[T1]{fontenc}
\usepackage[french]{babel}
\usepackage{amsmath}
\usepackage{amssymb}
\usepackage{booktabs}
\usepackage{array}
\usepackage{tikz}
\usetikzlibrary{positioning,calc,arrows.meta}
\usepackage{eurosym}

% Thème : celui des autres séances
\usetheme{Madrid}
\usecolortheme{default}

\definecolor{bleuactu}{RGB}{0,51,102}
\definecolor{vertactu}{RGB}{0,102,51}
\definecolor{rougeactu}{RGB}{153,0,0}
\definecolor{orangeactu}{RGB}{230,126,34}

\setbeamercolor{structure}{fg=bleuactu}
\setbeamercolor{frametitle}{bg=bleuactu,fg=white}

\newcommand{\keur}{\,k\euro}
% Un cas : cadre orange, renvoi à la séance qui le traite
\newenvironment{cas}[2]{\begin{exampleblock}{Cas #1 --- #2}}{\end{exampleblock}}

% Compilation pdfLaTeX sans erreur : caractères absents de l'encodage T1 (voir README « LaTeX »)
\DeclareUnicodeCharacter{2248}{\ensuremath{\approx}}
\DeclareUnicodeCharacter{2264}{\ensuremath{\leq}}
\DeclareUnicodeCharacter{2265}{\ensuremath{\geq}}

\title[Séance 1 -- Panorama]{Séance 1 : la comptabilité prudentielle en un coup d'œil}
\subtitle{Pourquoi, comment, quand, à quoi ça sert, et ses limites}
\author[F. Borel-Mathurin]{Fabrice BOREL-MATHURIN}
\institute[ACPR-DEAR et SAF Lyon 1]{
    ACPR - DEAR\\
    et\\
    Laboratoire SAF Université Claude Bernard Lyon 1
}
\date{2026-2027}

\begin{document}

\frame{\titlepage}

\begin{frame}{Plan de la séance (4 h)}
\tableofcontents
\end{frame}

\begin{frame}{Objectifs pédagogiques}
\begin{block}{À l'issue de cette séance, vous serez capable de :}
\begin{enumerate}
  \item Expliquer pourquoi les assureurs relèvent d'une comptabilité prudentielle
  \item Décrire en quelques phrases la mécanique de Solvabilité II : bilan économique, capital requis, ratio, trois piliers
  \item Situer les échéances : l'histoire du cadre, le calendrier d'une année, le moment où le superviseur intervient
  \item Dire à quoi sert le ratio de couverture, et ce qu'il ne dit pas
\end{enumerate}
\end{block}
\vspace{0.3cm}
\begin{alertblock}{Déroulé}
\textbf{Principes :} 1 h \quad | \quad \textbf{Cinq cas, en groupes :} 1 h 30 \quad | \quad
\textbf{Mise en commun :} 1 h \quad | \quad \textbf{La carte du cours :} 30 min
\end{alertblock}
\end{frame}

\begin{frame}{Pour commencer : trois questions}
\begin{block}{À discuter à deux, 5 minutes}
\begin{enumerate}
  \item Une obligation achetée 100\,\euro{} cote aujourd'hui 85\,\euro. Que vaut-elle dans les
        comptes de l'assureur ? Et pour quelqu'un qui voudrait la racheter ?
  \item Un assureur encaisse la prime \emph{avant} de payer les sinistres. Qu'est-ce que cela
        change pour celui qui le contrôle ?
  \item Qui a intérêt à ce qu'un assureur reste solvable ?
\end{enumerate}
\end{block}
\vspace{0.2cm}
\small Nous reviendrons sur vos réponses au fil de la séance.
\end{frame}

% =====================================================================
\section{Pourquoi ?}
% =====================================================================

\begin{frame}{Un métier à cycle inversé}
\begin{columns}[T]
\begin{column}{0.48\textwidth}
\begin{block}{L'assureur}
\begin{itemize}
  \item \textbf{encaisse d'abord} la prime, \textbf{paie ensuite}, parfois des décennies plus tard
  \item ne connaît son coût de revient qu'\textbf{après coup}
  \item porte des engagements \textbf{longs}, \textbf{aléatoires}, souvent assortis de \textbf{garanties}
\end{itemize}
\end{block}
\end{column}
\begin{column}{0.48\textwidth}
\begin{alertblock}{Le problème}
Le résultat de l'année ne dit pas si l'assureur pourra payer ses assurés \textbf{même après un
choc}. Il faut une mesure \textbf{prospective} et \textbf{économique} de sa solidité.
\end{alertblock}
\end{column}
\end{columns}
\vspace{0.4cm}
\begin{block}{Les objectifs du contrôle}
\textbf{Protéger les assurés et les bénéficiaires}, objectif principal (directive 2009/138/CE,
art.~27) ; préserver la \textbf{stabilité financière} ; maintenir la \textbf{confiance}.
\end{block}
\end{frame}

\begin{frame}{Des marges forfaitaires à une exigence fondée sur les risques}
\begin{columns}[T]
\begin{column}{0.48\textwidth}
\begin{block}{Solvabilité I (1973, 1979, 2002)}
Une marge \textbf{forfaitaire} : par exemple, en vie, 4\,\% des provisions mathématiques
$+$ 0,3\,\% des capitaux sous risque. Deux assureurs aux placements opposés ont la même exigence.
\end{block}
\end{column}
\begin{column}{0.48\textwidth}
\begin{block}{Solvabilité II (depuis 2016)}
Un \textbf{bilan économique}, et un capital requis qui dépend des \textbf{risques réellement
pris} : marché, crédit, souscription, opérationnel. La même logique que Bâle II pour les banques.
\end{block}
\end{column}
\end{columns}
\vspace{0.4cm}
\begin{cas}{A}{Equitable Life (Royaume-Uni, 2000)}
Des contrats de retraite vendus avec une option de rente à taux garanti, qui ne valait presque
rien à la vente. Les taux baissent : l'option devient très coûteuse, et rien ne l'avait mesurée.
\textbf{Question} : qu'aurait-il fallu valoriser, et qui paie à la fin ? \hfill
\emph{$\to$ séances 5 à 7}
\end{cas}
\end{frame}

% =====================================================================
\section{Comment ?}
% =====================================================================

\begin{frame}{La mécanique en une image}
\begin{center}
\begin{tikzpicture}[
  boite/.style={draw=bleuactu, fill=bleuactu!8, rounded corners, minimum width=2.9cm,
                minimum height=1.25cm, align=center, font=\small},
  fleche/.style={-{Latex[length=2mm]}, thick, bleuactu}]
\node[boite] (b) {\textbf{Bilan économique}\\\scriptsize valeur d'échange (art.~75)};
\node[boite, right=1cm of b] (f) {\textbf{Fonds propres}\\\scriptsize classés, plafonnés};
\node[boite, right=1cm of f] (r) {\textbf{Ratio}\\\scriptsize fonds propres / SCR};
\node[boite, below=1.1cm of f, draw=rougeactu, fill=rougeactu!8] (s) {\textbf{SCR}\\\scriptsize VaR 99,5\,\%, un an};
\draw[fleche] (b) -- node[above, font=\scriptsize]{actif $-$ passif} (f);
\draw[fleche] (f) -- (r);
\draw[fleche] (s) -| (r);
\end{tikzpicture}
\end{center}
\begin{itemize}
  \item \textbf{Actifs} en valeur de marché ; \textbf{provisions techniques} $=$ meilleure estimation $+$ marge de risque
  \item \textbf{SCR} : le capital qui permet d'encaisser un choc « une fois en 200 ans » sur un an
  \item \textbf{MCR} : un plancher, entre 25 et 45\,\% du SCR, sous lequel l'agrément est en jeu
\end{itemize}
\end{frame}

\begin{frame}{Trois piliers, deux autorités}
\begin{columns}[T]
\begin{column}{0.55\textwidth}
\small
\renewcommand{\arraystretch}{1.4}
\begin{tabular}{>{\raggedright\arraybackslash}p{2.4cm}>{\raggedright\arraybackslash}p{4.3cm}}
\toprule
\textbf{Pilier} & \textbf{Contenu} \\
\midrule
\textbf{1. Chiffre} & Bilan prudentiel, fonds propres, SCR, MCR \\
\textbf{2. Gouvernance} & Fonctions clés, gestion des risques, ORSA \\
\textbf{3. Information} & États au superviseur, rapport public (SFCR) \\
\bottomrule
\end{tabular}
\end{column}
\begin{column}{0.42\textwidth}
\begin{block}{EIOPA (européenne)}
\small Normes techniques, orientations, tests de résistance, courbe des taux sans risque mensuelle.
\end{block}
\begin{block}{ACPR (française)}
\small Agrément, contrôle sur pièces et sur place, sanctions, résolution.
\end{block}
\end{column}
\end{columns}
\vspace{0.3cm}
\begin{cas}{B}{l'Assureur Bêta}
Capitaux propres comptables : 50\,000\keur. Après revalorisation des actifs, des provisions et des
impôts différés : 23\,000\keur. SCR : 70\,000\keur. \textbf{Question} : où sont passés les
27\,000\keur, et l'assureur est-il couvert ? \hfill \emph{$\to$ séance 2}
\end{cas}
\end{frame}

% =====================================================================
\section{Quand ?}
% =====================================================================

\begin{frame}{Un cadre construit en cinquante ans}
\begin{center}
\begin{tikzpicture}[
  date/.style={font=\small\bfseries, text=bleuactu},
  evt/.style={font=\scriptsize, align=center, text width=2.1cm}]
\draw[-{Latex[length=3mm]}, very thick, bleuactu] (0,0) -- (14.2,0);
\foreach \x/\a/\t/\pos in {0.4/1973/{Directive non-vie : marge forfaitaire}/above,
                           2.4/1979/{Directive vie}/below,
                           4.4/2002/{Solvabilité I réformée}/above,
                           6.4/2009/{Directive Solvabilité II adoptée}/below,
                           8.4/2014/{Omnibus II : mesures de long terme}/above,
                           10.4/2016/{Entrée en application}/below,
                           12.6/2027/{Révision (directive (UE) 2025/2)}/above} {
  \fill[bleuactu] (\x,0) circle (2.5pt);
  \node[date, \pos=2pt] at (\x,0) {\a};
  \node[evt, \pos=16pt] at (\x,0) {\t};
}
\end{tikzpicture}
\end{center}
\vspace{0.2cm}
\small Chaque étape répond à une limite de la précédente : des marges sans lien avec le risque
(Solvabilité I), puis un bilan de marché trop volatil pour les engagements longs (mesures de
2014), puis des exigences jugées trop lourdes ou mal calibrées (révision).
\end{frame}

\begin{frame}{Le calendrier d'une année}
\begin{table}
\centering
\small
\renewcommand{\arraystretch}{1.35}
\begin{tabular}{p{4cm}p{9cm}}
\toprule
\textbf{Quand} & \textbf{Quoi} \\
\midrule
Chaque trimestre & MCR recalculé (art.~129) ; états trimestriels au superviseur \\
Au moins une fois par an & SCR calculé (art.~102) ; états annuels ; rapport public (SFCR) ; ORSA \\
Dès que le profil de risque change sensiblement & SCR recalculé, ORSA refait \\
En permanence & L'entreprise doit couvrir son SCR et son MCR, pas seulement aux dates de remise \\
\bottomrule
\end{tabular}
\end{table}
\begin{block}{Quand le superviseur intervient}
\small
SCR non couvert : \textbf{plan de rétablissement} sous deux mois, retour à la couverture en six
mois (art.~138). MCR non couvert : \textbf{plan de financement} sous un mois, retour en trois mois
(art.~139), puis, à défaut, \textbf{retrait d'agrément} (art.~144).
\end{block}
\end{frame}

\begin{frame}{Un cas pour « quand »}
\begin{cas}{C}{le cycle des taux, 2012--2022}
Un assureur vie a vendu des contrats d'épargne avec un taux minimum garanti de 2 à 3\,\%. De 2012
à 2021, les taux de marché baissent jusqu'à devenir négatifs sur certaines échéances ; en 2022,
ils remontent brutalement.
\begin{itemize}
  \item Taux en baisse : que deviennent la valeur des obligations, la meilleure estimation des
        contrats, le ratio ?
  \item Taux en hausse rapide : quel nouveau risque apparaît pour des contrats rachetables ?
  \item À quel rythme l'assureur et l'ACPR voient-ils ces effets ?
\end{itemize}
\hfill \emph{$\to$ séances 5 à 7}
\end{cas}
\end{frame}

% =====================================================================
\section{À quoi ça sert ?}
% =====================================================================

\begin{frame}{Un même chiffre, plusieurs usages}
\begin{table}
\centering
\small
\renewcommand{\arraystretch}{1.4}
\begin{tabular}{p{3.3cm}p{9.7cm}}
\toprule
\textbf{Pour qui} & \textbf{À quoi ça sert} \\
\midrule
Les assurés & Être payés, même après un choc sévère ; à défaut, une intervention précoce \\
Le superviseur & Détecter tôt, comparer les entreprises sur une base commune, graduer son intervention \\
La direction et le conseil & Piloter : dividendes, réassurance, allocation d'actifs, conception des produits, appétit au risque \\
Le marché et le public & Comparer, juger la solidité, grâce au rapport public annuel (SFCR) \\
\bottomrule
\end{tabular}
\end{table}
\vspace{0.2cm}
\begin{cas}{D}{l'Assureur Alpha}
Fonds propres éligibles 120\,M\euro, SCR 100\,M\euro : ratio de 120\,\%. Un choc de marché coûte
20\,M\euro. Le conseil devait voter un dividende de 10\,M\euro. \textbf{Question} : le peut-il ?
Quels leviers pour remonter le ratio ? \hfill \emph{$\to$ séances 8 et 9}
\end{cas}
\end{frame}

% =====================================================================
\section{Et ses limites}
% =====================================================================

\begin{frame}{Ce que le ratio ne dit pas : un aperçu}
\begin{table}
\centering
\small
\renewcommand{\arraystretch}{1.3}
\begin{tabular}{p{4.1cm}p{7.6cm}c}
\toprule
\textbf{Limite} & \textbf{En une phrase} & \textbf{Séance} \\
\midrule
Volatilité, procyclicité & Une valeur de marché bouge avec les marchés ; vendre en bas de cycle aggrave la crise & 6--8 \\
Une seule mesure & Un quantile à 99,5\,\% sur un an, pas une prévision ni le pire scénario & 8 \\
Un calibrage moyen & La formule standard est calibrée pour un assureur « type » & 8 \\
Modèles et données & Hypothèses, générateurs de scénarios, qualité des données : autant de sources d'erreur & 3--7 \\
Complexité et coût & Des calculs lourds, d'où le principe de proportionnalité & 2, 8 \\
Risques mal saisis & Liquidité, climat, cyber : traités surtout dans l'ORSA & 9 \\
\bottomrule
\end{tabular}
\end{table}
\end{frame}

\begin{frame}{Un cas pour les limites}
\begin{cas}{E}{le printemps 2020}
Mars 2020 : les actions chutent, les écarts de crédit s'élargissent, les taux restent très bas.
Deux mécanismes du cadre amortissent la chute des ratios : l'\textbf{ajustement symétrique} du
choc actions et la \textbf{correction pour volatilité} de la courbe d'actualisation. En avril,
l'EIOPA appelle les assureurs à suspendre leurs dividendes.
\begin{itemize}
  \item Pourquoi un bilan en valeur de marché peut-il amplifier une crise ?
  \item Que corrigent ces deux mécanismes, et à quel prix ?
  \item Après un choc record, le SCR calibré sur le passé mesure-t-il encore le bon risque ?
\end{itemize}
\hfill \emph{$\to$ séances 6 à 8}
\end{cas}
\vspace{0.1cm}
\small \emph{Pour aller plus loin} : la réévaluation du capital requis après un grand choc est
une question de recherche ; F. Borel-Mathurin et S. Loisel, \emph{Re-evaluation of the capital charge
in insurance after a large shock}, article soumis.
\end{frame}

% =====================================================================
\section{La carte du cours}
% =====================================================================

\begin{frame}{La carte du cours}
\begin{table}
\centering
\small
\renewcommand{\arraystretch}{1.25}
\begin{tabular}{clp{6.4cm}}
\toprule
\textbf{N°} & \textbf{Séance} & \textbf{Les cas qui y reviennent} \\
\midrule
2 & Le bilan prudentiel & B : l'Assureur Bêta, en entier \\
3 & Excel, Python et agents d'IA & Les contrôles d'un superviseur sur six assureurs \\
4 & Meilleure estimation non-vie & --- \\
5 & Meilleure estimation vie & A : la valeur d'une garantie ; C : rachats \\
6--7 & Approche stochastique et ALM & A, C, E : options, taux, adossement \\
8 & SCR, MCR et fonds propres & D : l'Assureur Alpha ; E : l'ajustement symétrique \\
9 & ORSA & D : appétit au risque ; E : risques hors formule \\
\bottomrule
\end{tabular}
\end{table}
\end{frame}

\begin{frame}{À retenir}
\begin{enumerate}
  \item \textbf{Pourquoi} : l'assureur promet avant de payer ; il faut mesurer, à valeur
        économique, s'il pourra tenir sa promesse après un choc.
  \item \textbf{Comment} : un bilan en valeur d'échange, des fonds propres classés, un SCR à
        99,5\,\% sur un an, un ratio ; trois piliers, deux autorités.
  \item \textbf{Quand} : en permanence, avec un MCR trimestriel, un SCR au moins annuel, et une
        intervention graduée dès que le SCR n'est plus couvert.
  \item \textbf{À quoi ça sert} : protéger les assurés, contrôler tôt, piloter l'entreprise,
        informer le public.
  \item \textbf{Ses limites} : volatilité, mesure unique, calibrage moyen, modèles : le ratio
        éclaire, il ne remplace pas le jugement.
\end{enumerate}
\end{frame}

\begin{frame}{Pour la prochaine séance}
\begin{block}{À faire}
\begin{itemize}
  \item Relire les cinq cas et vos réponses ; le corrigé est publié sur le site après la séance
  \item Parcourir les unités 1 et 2 du \textbf{polycopié de prérequis} : valeur actuelle, échéancier
\end{itemize}
\end{block}
\begin{block}{Questions de rappel, en ouverture de la séance 2}
\begin{itemize}
  \item Pourquoi le résultat comptable ne suffit-il pas à juger si un assureur pourra payer ses assurés ?
  \item Que mesure le SCR, et sur quel horizon ?
  \item Citez deux limites du ratio de couverture.
\end{itemize}
\end{block}
\end{frame}

\begin{frame}[c]
\centering
\Huge Questions ?

\vspace{1cm}
\Large
\textbf{Séance 2} : le bilan prudentiel
\end{frame}

\end{document}
